\section{Selecci\'on de bibliotecas y primitivas}
\label{sec:primitivas}

\begin{table}[ht]
\centering
\small
\begin{tabularx}{\textwidth}{@{}l l X@{}}
\toprule
Para qu\'e & Qu\'e usamos & Por qu\'e \\
\midrule
Grupo matem\'atico &
\textbf{ristretto255} (RFC~9496 \cite{rfc9496}) &
Grupo c\'iclico de orden primo $q \approx 2^{252}$. Elimina el cofactor 8 de
Curve25519, con lo que desaparecen los ataques de subgrupo peque\~no
\cite{chou-orlandi-simplestOT}. Cada elemento tiene una \'unica codificaci\'on
can\'onica de 32\,B y la decodificaci\'on rechaza las dem\'as. \\
Derivar claves &
\textbf{HKDF-SHA256} (RFC~5869 \cite{rfc5869-hkdf}) &
Convierte el secreto Diffie-Hellman $Z_j$ en una clave de 256 bits. La sal
$\mathsf{ctx}$ separa cat\'alogos y versiones; el \code{info} ata el dominio
\code{"oblivion-ti/ot-v1"}, el $sid$ y el \'indice $j$. \\
Cifrado sim\'etrico &
\textbf{AES-256-GCM} (NIST SP 800-38D \cite{nist-gcm}) &
Cifrado autenticado (AEAD) con tag de 128 bits y nonce aleatorio de 96 bits.
Cada clave ($K_j$ y $k_j$) cifra un solo mensaje, as\'i que no hay riesgo de
repetir nonce. El dato asociado ata contexto, sesi\'on e \'indice. \\
\bottomrule
\end{tabularx}
\caption{Primitivas criptogr\'aficas seleccionadas y su justificaci\'on.}
\label{tab:seleccion-primitivas}
\end{table}

\subsection{Bibliotecas}
AES-256-GCM y HKDF-SHA256 se toman de \code{cryptography} (sobre OpenSSL,
con aceleraci\'on por hardware). Para el grupo, el protocolo necesita tres
operaciones sobre elementos arbitrarios: producto ($u_j = u_{j-1}\cdot v$),
exponenciaci\'on ($u_j^\beta$, $v^{-i}$) y validaci\'on de codificaciones.
\code{cryptography} no las ofrece: su API de curvas el\'ipticas se limita a
ECDH y ECDSA y no permite sumar dos puntos. \code{PyNaCl} tampoco: se
comprob\'o en las versiones 1.5.0 y 1.6.2 que expone esa aritm\'etica para
Ed25519, cuyo grupo tiene cofactor 8, pero no las funciones de
\texttt{ristretto255}. Se adopta por eso \code{rbcl}, que empaqueta
\code{libsodium} y expone sus funciones \code{crypto\_core\_ristretto255\_*}
(producto, validaci\'on de puntos y aritm\'etica de escalares) y
\code{crypto\_scalarmult\_ristretto255} (exponenciaci\'on). Como el producto
no valida sus entradas, todo elemento recibido se comprueba antes con
\code{is\_valid\_point}. Como respaldo queda NIST P-256, tambi\'en de orden
primo, con \code{PyCryptodome}, cuya clase \code{EccPoint} permite sumar y
multiplicar puntos; solo cambiar\'ian \code{group\_id} y la serializaci\'on
de los elementos.

Las dependencias se fijan a versiones exactas en \code{requirements.txt},
sobre Python~3.11: rbcl 1.1.2, cryptography 50.0.2, fastapi 0.142.2, uvicorn
0.54.0, typer 0.27.2, pydantic 2.13.5 y httpx 0.28.1 (cliente HTTP).
